\section{Introduction}
\label{sec:intro}

Spontaneous parametric down-conversion (SPDC) remains the workhorse
process for generating photon pairs and, by detection of one photon of
a pair, heralded single photons; it underpinned the first two-photon
interference experiments \cite{hong1987} and the bright
polarization-entangled sources that followed \cite{kwiat1995}. For
modern applications the heralded photon must also be spectrally pure.
When the joint spectral amplitude (JSA) of the pair carries
signal--idler frequency correlations, the heralding measurement leaves
its partner in a mixed spectral state, and photons from independent
sources become partially distinguishable, degrading quantum
interference \cite{grice1997}. The Schmidt decomposition makes this
quantitative: the heralded purity is the inverse of the Schmidt number
of the JSA \cite{law2000}. Matching the group velocities of pump and
down-converted fields can render the JSA approximately factorable at
the source \cite{uren2005}, a strategy demonstrated experimentally with
pure heralded photons produced without lossy spectral filtering
\cite{mosley2008}.

A mature source-engineering literature has since pushed intrinsic
purity toward unity, most prominently by apodizing the nonlinear
response through domain engineering of the crystal
\cite{branczyk2011,graffitti2018}, and spectrally pure heralded sources
are now a standard requirement in single-photon-source design
\cite{meyerscott2020}. Joint-spectral characterization has also shown
that the pump spectrum itself -- not only its nominal bandwidth -- can
limit the purity a real source achieves \cite{zielnicki2018}.

In nearly all of this work, however, the pump laser is idealized as a
noiseless, transform-limited, single-longitudinal-mode Gaussian. Real
pump lasers deviate from this ideal in ways that are already
quantified, at no experimental cost, on the manufacturer's datasheet:
center-frequency drift and jitter, relative intensity noise (RIN), and
the longitudinal-mode structure of the laser cavity. What has been
missing is a translation between those datasheet lines and the figure
of merit a source builder cares about. How many gigahertz of pump
jitter cost one percent of heralded purity? Does RIN matter at all?

This paper provides that translation for a concrete, numerically
transparent model. We study a collinear SPDC source described by
first-order group-velocity-mismatch (GVM) phase matching, with
parameters representative of a 10 mm PPKTP crystal pumped at 775 nm and
emitting at 1550 nm, with a transform-limited Gaussian pump of 0.164
THz bandwidth (standard deviation); the ideal heralded purity of this
reference design is $P = 0.810$. Each
pump imperfection is mapped to a Monte-Carlo ensemble of pure pump
realizations, and the heralded photon is described by the
ensemble-averaged density matrix, from which we compute the mixed-state
purity directly. The results form a compact imperfection budget
(Table~\ref{tab:budget}). Center jitter of 26, 61, and 88 GHz rms
consumes a 1\%, 5\%, and 10\% relative purity budget, respectively;
re-matching the pump to a 5 mm crystal roughly doubles these
tolerances, at the cost of pair rate, which we do not model. Even a
single additional longitudinal mode of equal power at 0.25 THz spacing
costs 17\% of the purity -- more than the entire 10\% budget; the cost
scales with the side mode's relative power, so a weak, well-suppressed
side mode is far less damaging. RIN, by contrast,
leaves the heralded spectral purity exactly invariant in this model but
inflates accidental double-pair emission by the factor $1 + r^2$ for
fractional intensity noise $r$. Datasheet lines therefore map onto two
distinct damage channels: spectral imperfections that degrade purity,
and statistical imperfections that degrade pair statistics.

We emphasize the scope: this is a numerical modeling study with
first-order (linear) GVM phase matching, quasi-static ensemble noise,
and heralded spectral purity (together with the relative double-pair boost)
as the figures of merit; absolute pair rate, heralding efficiency, and
higher-order dispersion lie outside the model. Section~\ref{sec:model} defines the JSA, the
Schmidt-decomposition purity, and the mixed-state ensemble method.
Section~\ref{sec:results} establishes the ideal design baseline
(pump-bandwidth and crystal-length optima) and presents the
imperfection sweeps behind the budget. Section~\ref{sec:discussion}
summarizes the budget, develops the spectral-versus-statistical reading
of a laser datasheet, and discusses limitations and extensions.
